\section{One analytic band for all positive counts}\label{simp:scalar-band}

Keep the scalar state and localized transfer of
Sections~\ref{simp:scalar-state} and \ref{simp:controller-section}.
Set
\[
 t_-=1199/10000,\quad t_+=3001/10000,\quad d_*=23/20,
 \quad U_0=.6905799,\quad\lambda_0=.2319048039,\quad\xi=10^{-6}.
\]
The two ordinary constraints used here are
\[
 y_j+\tfrac25(t-S_j)\le R_4e^{\pi_j},\qquad
 y_j+\tfrac12(t-S_j)\le R_5e^{\pi_j}\quad(D\le4y_j),
 \qquad R_5=1.02631852.
\]
Write $\widehat P_j=P_j-\frac12\min(\pi_j,K)$ and
$F(H)=\mathcal E(H)+b+\Psi(H)$.
The target $b$ is the actual row target; the proof allows every
$b\in[0,b_c]$ and does not require equality at a clipped multiplier.

\begin{theorem}\label{simp:band}
If $t_-\le t\le t_+$ and $0\le D\le d_*$, then, for every
head set with $1\le h=|H|\le4$ and arbitrary nonnegative prices,
\begin{equation}\label{simp:band-result}
 F(H)<U_0+(h-1)\lambda_0-\xi.
\end{equation}
The proof includes its own one-head base and all unbounded heads.
Thus the same whole-box rule is independently valid for all four
numerical verifiers.
\end{theorem}

\subsection{A constant credit and one owner inequality}

Put $d=\max(1,D)$. Since $0\le v(r)\le r$, for $D\le1$
and $T\le t$ one has
\[
 D[1+v((T/D-\beta)^+)]\le\max(D,T+zD)\le1;
\]
here $t_++z<1$. For $D\ge1$, $t_+<\beta$ gives $W(T,D)=D$.
Consequently the credit of the controller is constant in $\mu$:
\begin{equation}\label{simp:band-credit}
 W(t-Q,D)=d,\qquad
 C(t,D,0;\mu)=c(D):=\min\{(3/5)\log d,h_0(\log d)\}.
\end{equation}
Indeed the minimizing $Q$ is zero and
$\log M(D)\ge\log(481/400)>(3/5)\log d_*$ when $D\ge1$.
The HOME constraints also give $\phi_H\ge c(D)\ge0$.
The function $c(D)+1+\log r-rt$ increases strictly on $(0,1]$,
starts at $-\infty$, and ends above $b_c$.
Moreover
\[
 k<21/50,\qquad
 1+\log(21/50)-(21/50)t_+>b_c.
\]
Therefore the actual clipped controller satisfies
\begin{equation}\label{simp:band-mu}
 k/d\le\mu\le21/50,
 \qquad c(D)-\mu t+1+\log\mu\ge b.
\end{equation}

For a head write $w=e^\pi$. In addition to $\widehat P\ge0$,
the price envelope gives
\begin{align}
 \widehat P&\ge\tfrac{11}{40}(w-1)quad(1\le w\le2),
                                                   \label{simp:band-price-chord}\\
 \widehat P&\ge\log w+\log y-\log(y+5Dw/16)-K/2,
                                                   \label{simp:band-parallel-price}\\
 e^{-\widehat P}&\le e^{K/2}(w^{-1}+5D/(16y))\quad(y>0).
                                                   \label{simp:band-tail-price}
\end{align}
For completeness, \eqref{simp:band-price-chord} follows from
$\widehat P\ge\max\{g(\pi),(2/5)\pi\}-\frac12\min(\pi,K)$.
On $[0,K]$, use $g(\pi)\ge\pi-5\pi^2$ and
$e^\pi\le1+\pi+(3/5)\pi^2$; the remaining coefficient is
\[
 \tfrac12-\tfrac{11}{40}-(5+\tfrac35\tfrac{11}{40})a_0>0.
\]
On $[K,a_0]$ the difference has positive derivative, since
$e^{-10a_0}>(11/40)e^{a_0}$.
On $[a_0,1/400]$ use $b_0-K/2>(11/40)(e^{1/400}-1)$.
On $[1/400,\log2]$ the difference
$(2/5)\pi-K/2-(11/40)(e^\pi-1)$ is concave with positive endpoints.
The last two price bounds follow from reservation and
$M(w\min(1,D/y))\le1+5Dw/(16y)$.

Define
\[
 q_0(t)=\min\{R_4-(2/5)t,R_5-t/2\},\quad
 q_-=R_5-t_+/2,\quad q_+=R_4-(2/5)t_-.
\]
For every head with $y\le6$ we claim
\begin{equation}\label{simp:band-owner}
 \mu\beta y-\mu zS-\widehat P\le\mu\beta q_0(t).
\end{equation}
Let $\Delta=y-S/2-q_0(t)$. If $\Delta\le0$, use
$z>\beta/2$ and $\widehat P\ge0$.
Otherwise $y>q_->4/5$, so both ordinary constraints are available.
They imply
$y-S/2\le q_0(t)+R_5(w-1)$, hence $w\ge1+\Delta/R_5$.
For $0<\Delta\le R_5$, \eqref{simp:band-price-chord} gives
$\widehat P\ge\mu\beta\Delta$, since
$(21/50)\beta R_5<11/40$.
When $\pi\ge\log2$, use instead
$(2/5)\log2-K/2>11/40$ to obtain the same bound.

For $R_5\le\Delta\le6-q_-$, let $c_M=23/64$ and
\[
 L_*(\Delta)=\log(1+\Delta/R_5)+\log(q_-+\Delta)
 -\log\{q_-+\Delta+c_M(1+\Delta/R_5)\}.
\]
The function $\log w+\log y-\log(y+c_Mw)$ increases in both
variables and is jointly concave. In relative differential coordinates
$A=dw/w,B=dy/y$, its Hessian form is
\[
 -A^2-B^2+(\lambda A+(1-\lambda)B)^2\le0,
 \qquad \lambda=c_Mw/(y+c_Mw)\in(0,1).
\]
Thus \eqref{simp:band-parallel-price} bounds $\widehat P$ below by
$L_*(\Delta)-K/2$. The concave function
$L_*(\Delta)-K/2-(21/50)\beta\Delta$ has margins greater than
$0.1000$ and $0.1233$ at the two endpoints.
It is positive throughout; $z>\beta/2$ now proves
\eqref{simp:band-owner} in every case.

\subsection{Two simultaneous scalar bounds}

Suppose first that all heads have $y_j\le6$.
The shared-overlap tangent uses
\[
 -\mu\beta(\sum_jS_j-\gamma)^+-\mu z\gamma
 \le-\mu z\sum_jS_j,
\]
which follows from $\beta>z$ by considering either sign of the
positive-part argument. Summing \eqref{simp:band-owner} gives two
bounds at the same actual multiplier:
\begin{align}
 F(H)&\le A_h(\mu;t):=-\log\beta+h\mu\beta q_0(t),
                                                      \label{simp:band-A}\\
 F(H)&\le B_h(\mu;t,c):=-\log(\mu\beta)-1+\mu t-c+b_c
                         +h\mu\beta q_0(t)quad(\phi_H\ge c\ge0).
                                                      \label{simp:band-B}
\end{align}
For the first bound use $\phi_H+\mu q\ge c(D)$ and
\eqref{simp:band-mu}; the actual $b$ cancels exactly.
For the second discard $q\ge0$ and use $b\le b_c$.

For $h\le3$, $B_h$ decreases throughout $0<\mu\le21/50$:
\[
 \partial_\mu B_h\le-50/21+t_++3\beta q_+<-0.2748.
\]
Let $r_c(t)$ solve $c+1+\log r_c-r_ct=b_c$ below $21/50$.
The two bounds agree at this root, while $A_h$ increases. Hence
\begin{equation}\label{simp:band-benchmark}
 \min(A_h,B_h)\le-\log\beta+h\beta r_c(t)q_0(t).
\end{equation}
This argument includes every clipped-floor state.
Set $t_c=10(R_5-R_4)=.2599495$.
On a segment $q_0=R-st$,
\[
 \frac{d}{dt}\{r_c(t)(R-st)\}
 =\frac{r_c(t)(r_c(t)R-s)}{1-r_c(t)t}.
\]
The root increases with $t$, so the first score segment has no
interior maximum; the second decreases because $(21/50)R_5<1/2$.

Three regions cover all bounded-head states for $h=1,2,3$.
For $t\ge3/20$, take $c=0$ and test $t=3/20,t_c$ in
\eqref{simp:band-benchmark}.
For $D\ge1003/1000$, take $c_*=(3/5)\log(1003/1000)$ and
test $t=t_-,t_c$. Every HOME branch exceeds $c_*$:
the localized branch uses $(2/5)\log(1003/1000)>b_0$, and
the others use monotonicity and $M(1)=481/400$.
The root upper bounds needed for these tests are
\[
 r_0(3/20)<.390121814249,\quad r_0(t_c)<.409249486646,
 \quad r_{c_*}(t_-)<.384621435921,
 \quad r_{c_*}(t_c)\le r_0(t_c).
\]
Each displayed rational endpoint has strictly positive defining
root residual, so monotonicity proves its direction.
Finally, on $t\le3/20,D\le1003/1000$,
$\mu\in[k/(1003/1000),k]$, since
$1+\log k-k(3/20)>b_c$.
Use \eqref{simp:band-B} with $c=0$, $\mu=k/(1003/1000)$,
and $t=3/20$. The last substitution increases the bound because
$1-(2h/5)\beta>0$ for $h\le3$.
For each of these finitely many endpoints the margin below
$U_0+(h-1)\lambda_0$ exceeds $1.7\cdot10^{-5}>\xi$.

\subsection{Four heads and the exact transition}

Put
\[
 t_f=(1+\log k-b_c)/k,
 \qquad .18<t_f<.183<t_c.
\]
If $t\le t_f$, then $\mu\in[k/d,k]$.
Every HOME branch dominates $G_D=\log d-1+1/d$ on $D\le d_*$.
For the linear branch the difference has derivative
$(1-(2/5)D)/D^2\ge0$; for the multiplicative branch use
$\log(481/400)>G_{d_*}$.
For the localized branch, put $s=\log D\ge0$.
If $s\le a_0$,
\[
 h_0(s)\ge(5-(50/3)a_0)s^2\ge s^2/2\ge s-1+e^{-s}.
\]
If $s\ge a_0$, use $1-e^{-a_0}>b_0$.
The $D\le1$ case is nonnegativity.

The right side of \eqref{simp:band-B}, with $h=4,c=G_D$, is
convex in $\mu$. Since $q_0=R_4-(2/5)t$, its two endpoints
$k/d,k$ yield
\begin{equation}\label{simp:band-four-floor}
 F(H)\le\log4-
 \frac{(2/5-k)t-(R_4-1)}d+b_c.
\end{equation}
At $k/d$, the HOME term cancels $\log d-1+1/d$ exactly;
at $k$, use $G_D\ge0$ and the positive numerator.
The largest bound is at $t=t_-,d=d_*$.
Its margin below $\log4$ exceeds $2.9323763\cdot10^{-6}$,
and below $U_0+3\lambda_0$ exceeds $2.88295\cdot10^{-6}>\xi$.

If $t\ge t_f$, then $r_0(t)\ge k$ and $c(D)\ge0$ imply
$\mu\le r_0(t)$ for every actual target $b\le b_c$.
Use \eqref{simp:band-A}; the preceding derivative calculation reduces
the maximum to $t_f$ or $t_c$. At $t_f$ the root equals $k$,
and at $t_c$ use its displayed rational upper bound.
Both endpoint margins exceed $1/50000$.
This proves the theorem for all bounded-head states.

\subsection{A direct one-head base closes every tail}

For one head with $y\ge6$, keep the $a$-penalty and discard
only nonnegative penalties on $\gamma,q,\phi_H$.
Since $V_+\le a+y$ and
$\mu\beta\ge1/(4d_*)>1/6\ge1/y$,
the maximum of $\log(a+y)-\mu\beta a$ occurs at $a=0$.
Using \eqref{simp:band-tail-price} and $w\ge y$ gives
\[
 F(\{j\})\le\log y-\widehat P+\mu t+b
 \le\log(1+5d_*/16)+K/2+(21/50)t_++b_c
 <.434<U_0-\xi.
\]
Thus the one-head result covers every possible head size.

For a set with $h\ge2$, nonnegative submodularity gives
$\sum_jV_+(H\setminus\{j\})\ge(h-1)V_+(H)$.
If some inverse charged price obeys
$e^{-\widehat P_j}\le(h-1)(e^\delta-1)$, all others are at
most one, so their sum is at most $(h-1)e^\delta$.
Exponentiating and summing the proposed deletion bounds then shows
that some deletion increases the scalar by at most $\delta$.
This is an elementary value-and-price statement, independent of any
positive-count theorem outside this section.

For every head with $y\ge6$, \eqref{simp:band-tail-price} gives
\[
 e^{-\widehat P}\le\frac{29}{128}e^{K/2}
                  <e^{\lambda_0-\xi}-1.
\]
Apply the deletion with $\delta=\lambda_0-\xi$ and retain the same
$a,\gamma,q,t,D,\mu,b,\phi_H$ and remaining prices.
Starting from the just-proved one-head result gives the two-head,
then the three-head, then the four-head result, with at least the
claimed reserve at every step. If no large head exists, the direct
bounded-head proof applies. This completes Theorem~\ref{simp:band}.

All signs above are checked by a single directed endpoint program,
including the three root brackets and each count's endpoint bound.
A numerical box wholly contained in this rectangle may therefore
be accepted with reserve $10^{-6}$ for any $h=1,2,3,4$.
